\section{Project Goal}
The primary goal is straightforward to state but less straightforward to answer cleanly: is Norwegian a more effective transfer source for Danish NER than English? The hypothesis behind it is reasonable — languages sharing closer history and structure should induce more compatible representations inside a multilingual encoder, which should translate into stronger cross-lingual transfer. Rather than treating this as self-evident, this study operationalizes the intuition as a controlled empirical comparison.

The practical angle matters as much as the theoretical one. Raw F1 is useful, but the real decision facing a practitioner is: when Danish annotation is scarce, which source language and encoder combination is the best bet? A consistent Norwegian advantage makes the choice obvious. A negligible difference shifts the decision toward corpus size and encoder selection instead.

A secondary aim is to untangle three influences that tend to get conflated: source-language choice, encoder architecture, and the amount of available Danish supervision. Placing monolingual baselines, zero-shot transfer, fine-tuned transfer, and Danish-only models side by side in a single study makes it possible to assign credit to each factor rather than leaving them entangled.
